\documentclass[10pt,a4paper]{amsart}
\usepackage[margin=19mm]{geometry}
\usepackage{amsmath,amssymb,mathtools}
\usepackage[scr=rsfs,cal=euler]{mathalfa}
\usepackage{xfrac}
\usepackage[unicode,hidelinks,bookmarks=false]{hyperref}
\newcommand{\ie}{i.e.\ }
\newcommand{\cf}{cf.\ }
\newcommand{\resp}{respectively\ }
\newcommand{\Subsection}{\subsection}
\providecommand{\nobreakdash}{\nobreak\hbox{-}}
\setlength{\emergencystretch}{2em}
\multlinegap0pt
\usepackage{bm}
\usepackage{bbm}
\usepackage{dsfont}
\usepackage{xspace}
\usepackage{cancel}
\usepackage{mathtools}
\usepackage{amsthm}
\makeatletter
\@ifundefined{equation}{\newtheorem{equation}{Equation}}{}
\@ifundefined{lemma}{\newtheorem{lemma}[equation]{Lemma}}{}
\@ifundefined{proposition}{\newtheorem{proposition}[equation]{Proposition}}{}
\makeatother
\DeclareMathOperator{\diam}{diam}
\newcommand{\N}{\mathbb{N}}
\newcommand{\ba}{\mathbf{a}}
\newcommand{\bbb}{\mathbf{b}}
\newcommand{\bc}{\mathbf{c}}
\title[Connecting conformal dimension and Poincar\'e profiles]{Connecting conformal dimension and Poincar\'e profiles}
\author{David Hume, John M. Mackay}
\date{Source: arXiv:2511.10469v1 (CC BY 4.0)}
\begin{document}
\maketitle

\noindent\emph{Source and licence.} Connecting conformal dimension and Poincar\'e profiles. Version arXiv:2511.10469v1 (CC BY 4.0), distributed under the Creative Commons Attribution 4.0 International licence, as recorded in the arXiv licence metadata for this exact version and shown on the abstract page https://arxiv.org/abs/2511.10469v1. Full rights notice: licenses/arxiv-2511.10469-CC-BY-4.0.txt.\par\smallskip

\noindent\textit{Editorial scope.} Counted unit: the Proposition whose source label is \texttt{prop:QI-round-tree-graph} in the article (source0.tex lines 902--903, source label \texttt{prop:QI-round-tree-graph}), delivered together with the article's own complete original proof of it (source0.tex lines 991--1008, one \texttt{proof} environment of 18 lines). No mathematics is written, repaired or re-derived: statement and proof are the article's own text line for line. The proof is detached from the statement in the source: the article states the result at lines 902--903 and prints its proof at lines 991--1008, and the proof environment's own header names the statement, so the association is the article's own. Required context retained uncounted: lem:cantorarcs0 (lines 914--916); lem:cantorarcs1 (lines 941--942); lem:cantorarcs3 (lines 971--973). Every definition, display and label of those lines is retained unchanged. Omitted from this package and not redistributed: 1--901, 904--913, 917--940, 943--970, 974--990, 1009--1289. Nothing inside a retained excerpt is omitted or altered: every retained body is delivered exactly as the article's lines are written, so no omission receipt of any kind accompanies this package. Citations: the retained excerpts carry no citation command at all, so no bibliography entry of the article is transcribed and no reference is redistributed. Reference adaptation: none was needed and none was made; every reference inside the delivered unit resolves inside it, so no unresolved reference or citation is produced. Printed slips kept literal: no printed word, symbol, hypothesis or number of the counted statement or of its proof was altered. Layout and class: the article's own class is \texttt{amsart} with packages including amsmath, amsthm, amsfonts, graphicx, amssymb, hyperref, xcolor, tikz; the wrapper is the lane's pinned \texttt{amsart} editorial wrapper at 10pt on A4, which restates the theorem-like environments the retained text needs and adds the article's own macro definitions that the retained text needs (\texttt{\textbackslash N}, \texttt{\textbackslash ba}, \texttt{\textbackslash bbb}, \texttt{\textbackslash bc}, \texttt{\textbackslash diam}), with extra wrapper packages (bm, bbm, dsfont, xspace, cancel, mathtools). The delivered numbering is the unit's own. Assets: no retained span contains a figure environment, a graphics-inclusion command, a PDF-page inclusion, a table or a TikZ picture; no image file is copied into this package, the package declares no asset dependency and no image file is redistributed with it. Licence: version arXiv:2511.10469v1 of this article is distributed under the Creative Commons Attribution 4.0 International licence, as recorded in the arXiv licence metadata for this exact version and shown on the abstract page https://arxiv.org/abs/2511.10469v1. A full rights notice is delivered at licenses/arxiv-2511.10469-CC-BY-4.0.txt. Characters: every retained excerpt is pure ASCII, so the delivered engine, XeLaTeX with its default Latin Modern text font, reports no missing character.
	\begin{lemma}\label{lem:cantorarcs0}
		The sets $W$ and $\{x_\ba^i: \ba \in \{0,1\}^*, 1 \leq i \leq N_{\ba}\}$ are at bounded Hausdorff distance.
	\end{lemma}

	\begin{lemma}\label{lem:cantorarcs1} For all non-empty $\ba \in \{0,1\}^*$ and $1 \leq i< N_\ba$, we have $\diam J_\ba[z_\ba^i, z_\ba^{i+1}] \in [\beta^n/4\lambda,\beta^n/2]$.  In particular, $\rho(z_\ba^i, z_\ba^{i+1}) \in [\beta^n/4\lambda,\beta^n/2]$.
	\end{lemma}

	\begin{lemma}\label{lem:cantorarcs3}
		Assume $\beta < 1/12 \lambda^2$.  Then there exists $H\geq 2$ so that for any $(a_1,\ldots,a_n)\in \{0,1\}^*$ and any $1 \leq j < N_{a_1\cdots a_{n-1}}$, if $i$ is minimal with $\pi_{a_1\cdots a_n}(i)=j$ and $i'$ is minimal with $\pi_{a_1\cdots a_n}(i')=j+1$, then $i+2 \leq i' \leq i+H$.
	\end{lemma}

	\begin{proposition}\label{prop:QI-round-tree-graph} There is an (irregular) round tree graph $A$ which is quasi-isometric to $W$.
    \end{proposition}

		\begin{proof}[Proof of Proposition \ref{prop:QI-round-tree-graph}]
		Recall $\{x_\ba^i : \ba \in \{0,1\}^*, 1 \leq i \leq N_\ba\}$ is coarsely dense in $W$ by Lemma \ref{lem:cantorarcs0}.
		The function $\Pi$ is a bijection onto its image, so it suffices to show that $\Pi$ and $\Pi^{-1}$ (restricted to its image) are coarsely Lipschitz.

		To show $\Pi$ is coarsely Lipschitz, since $A$ is a graph it suffices to bound the images of endpoints of edges.
		For a horizontal edge between $p_\ba^i$ and $p_\ba^{i+1}$ for some $\ba \in \{0,1\}^*, 1 \leq i < N_\ba$, by Lemma~\ref{lem:cantorarcs1} we have $\rho(z_\ba^i, z_\ba^{i+1}) \leq \beta^n/2$, hence the Gromov product $(z_\ba^i | z_\ba^{i+1})_1 \geq \frac{-1}{\epsilon}\log(\beta^n)-C$ for some $C$.  As each $x^j_\ba$ is a point $-\frac1\varepsilon\log(\beta^n)+C_1$ from $1$ along a geodesic to $z^j_\ba$, we deduce that $d(x_\ba^i, x_\ba^{i+1})$ is uniformly bounded by hyperbolicity.
		For a vertical edge between $p_\ba^k$ and $p_{\ba'}^j$, with $\ba=(\ba',a_n) \in \{0,1\}^n$, $1 \leq i\leq k \leq i'-1 \leq N_\ba$ and $i,i'$ minimal with $\pi_\ba(i)=j$ and $\pi_\ba(i')=j+1$ respectively, we have by Lemma~\ref{lem:cantorarcs3} (and proof) that $\rho(z_\ba^k, z_{\ba'}^j) \leq \beta^{n-1}$, so again the Gromov product satisfies $(z_\ba^k| z_{\ba'}^j)_1 \geq \frac{-1}{\epsilon}\log(\beta^n)-C'$, and $d(x_\ba^k, x_{\ba'}^j)$ is uniformly bounded by hyperbolicity.


		To show $\Pi^{-1}$ is coarsely Lipschitz, since $W$ is quasi-convex in $G$, $W$ is coarsely geodesic, and so it suffices to show that for any $M \geq 0$ there exists $M' < \infty$ so that $d(x_\ba^i, x_{\bbb}^j) \leq M$ implies $d_A(p_{\ba}^i,p_{\bbb}^j) \leq M'$.
		Assume we have such $x_\ba^i, x_\bbb^j$, with $\ba \in \{0,1\}^m, \bbb \in \{0,1\}^n$.
		Then $|n-m|$ is uniformly bounded depending on $M$; note that $\rho(z_\ba^i, z_\bbb^j) \preceq \beta^{m} \asymp \beta^{n}$.
		By Lemma~\ref{lem:cantorarcs0}, there exist $\hat\ba = (\ba, a_{m+1},\ldots) , \hat\bbb = (\bbb, b_{n+1},\ldots)\in \{0,1\}^\N$ with $\rho(z_\ba^i, J_{\hat\ba}) \preceq \beta^{m}$ and $\rho(z_\bbb^j, J_{\hat\bbb}) \preceq \beta^{n}$.
		Thus the minimum distance $\rho(J_{\hat\ba}, J_{\hat\bbb}) \preceq \beta^m$,
		and so if $k = \min\{l:a_l\neq b_l\}$, we have $\beta^k \preceq \beta^m$, which implies that $k \geq m-C$ for some constant $C$.  Hence for $\bc=(a_1,\ldots,a_k)=(b_1,\ldots,b_k)$, there exists $i', j'$ so that $d_A(p_\ba^i,p_\bc^{i'}), d_A(p_\bbb^j,p_\bc^{j'}) \leq C$ and $\rho(z_\ba^i,z_\bc^{i'}), \rho(z_\bbb^j, z_\bc^{j'}) \leq 2 \beta^{k+1}$.  Then $\rho(z_\bc^{i'}, z_\bc^{j'}) \preceq 4\beta^{k+1}+\beta^m \preceq \beta^k$.
		Using that $J_\bc$ is a $\lambda'$-quasi-arc, the chain of $\beta^k/4\lambda$-separated points $\{z_\bc^{k}: k \in [i',j']\} \subset B(z_\bc^{i'},C\lambda'\beta^k)$, and so $|j'-i'|$ is uniformly bounded by doubling.
		Hence $d_A(p_\bc^{i'},p_\bc^{j'})$ is uniformly bounded, and we are done.
	\end{proof}

\end{document}
